\documentclass{article}

\usepackage{amsmath}
\usepackage{amssymb}
\usepackage{amsthm}
\usepackage{geometry}
\usepackage{hyperref}
\usepackage{booktabs}
\usepackage{array}

\geometry{margin=1in}
\hypersetup{
  colorlinks=true,
  linkcolor=blue,
  urlcolor=blue
}

%% --- Project macros (shared with sts-paper.tex) ---------------------------
\newcommand{\STF}{{\normalfont\textsc{STF}}}
\newcommand{\SIR}{{\normalfont\textsc{SIR}}}
\newcommand{\STS}{{\normalfont\textsc{STS}}}

%% --- Paper-specific macros ------------------------------------------------
\newcommand{\Univ}{\mathcal{U}}
\newcommand{\CohTriple}{\mathrm{Coh}}
\newcommand{\ErrTriple}{\mathrm{Err}}
\newcommand{\ICE}{\mathrm{ICE}}
\newcommand{\Inf}{\mathrm{Inf}}
\newcommand{\Horizon}{\mathcal{H}}
\newcommand{\MA}{M_A}
\newcommand{\bottom}{\bot}

%% --- Theorem environments -------------------------------------------------
\newtheorem{theorem}{Theorem}
\newtheorem{corollary}[theorem]{Corollary}
\newtheorem{definition}{Definition}
\newtheorem{remark}{Remark}

%% ---------------------------------------------------------------------------
\title{Coherence as a Primitive:\\
A Formal Framework for Information, Cognition, and Computational Limits}
\author{Rogerio Figueiredo}
\date{2026}

\begin{document}

\maketitle

\begin{quote}
\textbf{Não normativo.}
Este documento apresenta contexto filosófico, hipóteses fundacionais e
motivação teórica para o sistema \STF{}-\SIR{}.
Ele \emph{não} define requisitos de implementação, contratos de API nem
critérios de aceitação de testes.
O posicionamento normativo do sistema é estabelecido em
\texttt{docs/adr/ADR-SEM-001-positioning.md}.
\end{quote}

\begin{abstract}
We propose a non-standard foundational hypothesis: that coherence, rather than
information, is the primary structural property underlying observable reality.
Motivated by an alternative interpretation of light as a static structural
dimension rather than a propagating entity, we formalize coherence as a
three-level relational invariant over proposition sets, distinguishing logical,
computational, and operational coherence.  We define \emph{Coherently Executable
Information} (CEI) as the subset of information that is both integrable into an
agent's existing model and capable of generating actionable inference, and prove
that information is usable if and only if it is coherently executable.  We show
that the inferability horizon of a finite agent corresponds precisely to the
complexity class~$\mathbf{P}$, establishing a formal correspondence between
epistemic limits and the $\mathbf{P}$ vs $\mathbf{NP}$ boundary.  We introduce
lexical coherence as a structure-preserving homomorphism between domains,
operationalized through the existing retention metric~$\rho$ of the \STF{}-\SIR{}
system, and classify lexical failures into four named error modes.  Finally, we
present a uniform coherence-error taxonomy---anomaly, hallucination, and
contradiction---each characterized by a three-component grounding vector.  The
resulting framework provides formal foundations for semantic compilation systems
such as \STF{}-\SIR{}, where coherence becomes a measurable, auditable invariant
of every knowledge transformation.
\end{abstract}

%% ---------------------------------------------------------------------------
\section{Introduction}

Classical information theory, following Shannon~\cite{shannon1948}, defines
information as the reduction of uncertainty in a random experiment.  This
account is remarkably effective for quantifying transmission fidelity, but it is
structurally agnostic: it assigns equal informational weight to messages that are
meaningful and to messages that are incoherent, provided they are equally
surprising.  A complementary theory of semantic information was developed by
Carnap and Bar-Hillel~\cite{carnap1953}, but even this account does not address
the question of whether a message can be \emph{integrated} into a receiving
system without destroying its existing structure.

We begin from a different starting point.  Rather than asking how much
uncertainty a message reduces, we ask: under what conditions is a message usable
at all?  Our answer is that usability requires coherence---and therefore that
coherence is logically prior to information.

This paper makes the following contributions.
\begin{enumerate}
  \item We introduce a foundational hypothesis ($H_0$) reinterpreting light as a
        static structural substrate rather than a propagating phenomenon, and
        derive from it the thesis that coherence is the primary property of the
        observable universe (\S\ref{sec:hypothesis}).
  \item We formalize coherence as a triple $(C_l, C_c, C_o)$ covering logical,
        computational, and operational dimensions, and characterize the eight
        states of the resulting coherence space (\S\ref{sec:coherence}).
  \item We define Coherently Executable Information (CEI) and prove that
        information is usable if and only if it is coherently executable
        (\S\ref{sec:information}).
  \item We establish the inference horizon of a finite agent as the set of
        propositions provable by polynomial-time algorithms, identifying
        epistemic limits with complexity-theoretic limits (\S\ref{sec:limits}).
  \item We reformulate perception as a structural matching (SAT) problem, in
        contrast to the standard signal-reception model (\S\ref{sec:perception}).
  \item We introduce lexical coherence as a coherence-preserving map between
        domains, operationalize it via the retention score~$\rho$ of \STF{}-\SIR{},
        and classify four modes of lexical failure (\S\ref{sec:lexical}).
  \item We present a uniform coherence-error taxonomy with three named classes
        and distinct detection strategies for each (\S\ref{sec:errors}).
  \item We connect the framework to \STF{}-\SIR{} by extending the semantic
        token model with derived coherence and information coordinates
        (\S\ref{sec:stfsir}).
\end{enumerate}

%% ---------------------------------------------------------------------------
\section{Foundational Hypothesis}
\label{sec:hypothesis}

\subsection{Static Light Substrate}

\begin{definition}[$H_0$: Static Light Substrate]
Light is not a propagating entity in spacetime.  It is a structural dimension
--- a static substrate --- over which sentient entities construct their
representation of the observable universe.
\end{definition}

This hypothesis inverts the standard kinematic model.  Rather than asking
\emph{how does light reach us?}, it asks \emph{what is light before it is
perceived as motion?}  Human perception of motion is inferential: no sensory
receptor directly captures velocity; what is captured are state differences
across successive temporal windows.  The interpretation of light as a traveling
entity is therefore a second-order cognitive construction, not a primary
observation.

Under $H_0$, light plays a role analogous to Euclidean space in classical
geometry: it is the \emph{medium} within which relations are defined, not an
entity moving within it.  This position is consonant with block-universe
interpretations of general relativity~\cite{minkowski1908}, in which past,
present, and future coexist as a static four-dimensional structure and
``motion'' is an artifact of an observer's worldline.

\subsection{The Primacy of Coherence}

The central ontological consequence of $H_0$ is as follows.

\begin{definition}[Thesis O: Coherence Precedes Information]
Information is not an ontological primitive.  It is derived from structural
compatibility.  Formally:
\[
  \text{information} \;\triangleq\; \text{that which does not break coherence.}
\]
\end{definition}

Under Thesis~O, no informational process can occur in the absence of an
underlying coherent structure.  Coherence is the condition of possibility for
any signal, message, or inference.

%% ---------------------------------------------------------------------------
\section{Formalizing Coherence}
\label{sec:coherence}

Let $\Univ$ be a universe of propositions and $\mathcal{R}$ a family of
consistency relations over $\Univ$.

\subsection{Three Levels of Coherence}

\begin{definition}[Logical Coherence, $C_l$]
A set $S \subseteq \Univ$ is \emph{logically coherent} if and only if
\[
  C_l(S) = 1 \;\iff\; \forall\, p, q \in S,\quad \mathcal{R}(p,q) \neq \bot,
\]
where $\bot$ denotes irreducible incompatibility.
\end{definition}

Logical coherence is a necessary condition for any useful system.  It does not,
however, guarantee that the system is tractable or productive.

\begin{definition}[Computational Coherence, $C_c$]
A set $S$ is \emph{computationally coherent} if its consistency can be decided
within a bounded computational effort.  The complexity-theoretic ideal is a
polynomial-time witness ($\exists\,\alpha \in \mathbf{P}$); in STF-SIR the
\emph{operational definition} is a step-budget proxy:
\[
  C_c(S) = 1 \;\iff\; \text{the coherence check terminates within }
  \mathit{step\_budget} \text{ steps}.
\]
The $\mathbf{P}$-membership framing (\S\ref{app:comp-coh}) provides theoretical
motivation; the step-budget threshold is the normative implementation criterion.
\end{definition}

\begin{remark}[Non-normative status of the P/NP framing]
The complexity-class characterization ($C_c = 0$ placing the problem in
$\mathbf{NP} \setminus \mathbf{P}$) is a theoretical bound, not an
implementation contract.  STF-SIR does \emph{not} claim to decide
$\mathbf{P} \neq \mathbf{NP}$; it approximates tractability via an explicit
\texttt{step\_budget} parameter.  Exceeding the budget sets $C_c = \mathit{Violated}$;
terminating within it sets $C_c = \mathit{Satisfied}$; the distinction is
operational, not complexity-theoretic.
\end{remark}

A set may be logically coherent yet computationally intractable within
a given step budget.  The table below treats $C_c$ as the result of the
step-budget check.

\begin{definition}[Operational Coherence, $C_o$]
A set $S$ is \emph{operationally coherent} if it supports productive inference:
\[
  C_o(S) = 1 \;\iff\; \exists\, f : S \to O \text{ with } O \neq \emptyset.
\]
\end{definition}

Operational coherence eliminates two pathologies invisible to $C_l$ and $C_c$:
tautological systems (logically valid, computationally verifiable, but
informationally inert) and systems that produce no actionable consequences.

\subsection{The Coherence Triple}

\begin{definition}[Coherence Triple]
The coherence of any set $S$ is fully characterized by
\[
  \CohTriple(S) = (C_l,\; C_c,\; C_o) \;\in\; \{0,1\}^3.
\]
\end{definition}

The eight elements of $\{0,1\}^3$ partition the space of possible coherence
states as follows.

\begin{center}
\renewcommand{\arraystretch}{1.25}
\begin{tabular}{cccl}
\toprule
$C_l$ & $C_c$ & $C_o$ & Classification \\
\midrule
0 & --- & --- & \textbf{Contradictory} --- invalid at base level \\
1 & 0   & --- & \textbf{Budget-exceeded} --- coherent but step-budget exhausted \\
1 & 1   & 0   & \textbf{Sterile} --- coherent, verifiable, but inert \\
1 & 1   & 1   & \textbf{Full coherence} --- the only productive state \\
\bottomrule
\end{tabular}
\end{center}

\begin{remark}
Coherence is not truth.  A system can satisfy $\CohTriple(S) = (1,1,1)$ and
be entirely disconnected from reality.  Coherence is an \emph{internal}
property; truth is an \emph{external} relation between model and world.
\end{remark}

%% ---------------------------------------------------------------------------
\section{Information as a Derived Concept}
\label{sec:information}

\subsection{Three Layers of Information}

\begin{definition}[Shannon Information, $I_1$]
For a discrete random variable $X$ with distribution $P$,
\[
  H(X) = -\sum_{x} P(x) \log_2 P(x).
\]
\end{definition}

Shannon entropy captures surprise and uncertainty but is blind to meaning.

\begin{definition}[Semantic Information, $I_2$ {\normalfont\cite{carnap1953}}]
The semantic information of a proposition $p$ in a possibility space
$\mathcal{W}$ is
\[
  \Inf(p) = -\log_2 P(\mathcal{W}_p),
\]
where $\mathcal{W}_p \subseteq \mathcal{W}$ is the set of worlds in which $p$
is true.
\end{definition}

Tautologies carry zero semantic information; contradictions carry maximal
semantic information but zero truth value.  Useful semantic information exists
only at the non-trivial intersection of coherence and truth.

\begin{definition}[Coherently Executable Information, $I_3$ / CEI]
\label{def:cei}
A message $m$ carries \emph{coherently executable information} for an agent $A$
with model $\MA$ if and only if
\[
  \ICE(m, A) = 1 \;\iff\;
  \begin{cases}
    \MA \cup \{m\} \not\models \bot
      & \quad \textbf{(C1: integrable)} \\
    \exists\, f : \MA \cup \{m\} \to O,\; O \neq \emptyset
      & \quad \textbf{(C2: executable)}
  \end{cases}
\]
\end{definition}

Condition C1 ensures the message does not destroy the existing model.
Condition C2 ensures the augmented model can derive at least one new
consequence.

\begin{theorem}[Information Usability Constraint]
\label{thm:usability}
A message is usable by agent $A$ if and only if it is coherently executable,
i.e., $\ICE(m,A) = 1$.
\end{theorem}

\begin{proof}
($\Rightarrow$) If $\MA \cup \{m\} \models \bot$, the model is contradicted and
no inference is valid.  If no $f$ exists such that $\MA \cup \{m\} \to O$ with
$O \neq \emptyset$, the message produces no actionable consequence.  In either
case, the message is unusable.

($\Leftarrow$) If both C1 and C2 hold, the message integrates without
contradiction and generates at least one new consequence.  It is therefore
usable.~$\square$
\end{proof}

\begin{corollary}
Any message satisfying C1 but not C2 is \emph{coherent noise}: it does not
invalidate the model but contributes nothing.  Any message satisfying C2 but not
C1 is a \emph{hallucination trigger}: it produces outputs but corrupts the base
model.
\end{corollary}

%% ---------------------------------------------------------------------------
\section{Cognitive and Computational Limits}
\label{sec:limits}

\subsection{The Inference Horizon}

Let $A$ be an agent with model $\MA$ and observation set $O_A$.

\begin{definition}[Inference Horizon]
The \emph{inference horizon} of agent $A$ is
\[
  \Horizon_A = \bigl\{\, p \in \Univ \;\mid\;
    \exists\, \alpha \in \mathbf{P},\;\; \alpha(\MA, O_A) \vdash p \,\bigr\}.
\]
\end{definition}

Everything outside $\Horizon_A$ is inferentially inaccessible to $A$, regardless
of computational resources applied.

\begin{theorem}[Computational Epistemic Boundary]
\label{thm:horizon}
The inferible knowledge of a finite agent is bounded by $\mathbf{P}$: the agent
can know exactly those propositions derivable by polynomial-time algorithms from
its model and observations.
\end{theorem}

\begin{proof}
By definition of $\Horizon_A$, any $p \in \Horizon_A$ is derived by some
$\alpha \in \mathbf{P}$.  For $p \notin \Horizon_A$, no polynomial-time
algorithm suffices; the agent may verify a proposed proof in $\mathbf{NP}$ but
cannot construct one efficiently.~$\square$
\end{proof}

\subsection{Three Epistemic Regimes}

\begin{center}
\renewcommand{\arraystretch}{1.25}
\begin{tabular}{ll}
\toprule
Regime & Epistemic status \\
\midrule
$\Pi \in \mathbf{P}$ &
  Computable and verifiable: fully within the agent's horizon \\
$\Pi \in \mathbf{NP} \setminus \mathbf{P}$ &
  Verifiable but not constructible: horizon of verification \\
$\Pi \notin \mathbf{NP}$ &
  Neither constructible nor verifiable: computationally unknowable \\
\bottomrule
\end{tabular}
\end{center}

\begin{corollary}[Coherence as Polynomial Witness]
Checking whether $\MA \models o$ for a single observation $o$ is polynomial in
$|\MA|$ for well-structured models.  Coherence verification is therefore within
$\mathbf{P}$, which explains the empirical preference of sentient agents for
coherent representations: coherence is the only regime in which verification is
computationally accessible.
\end{corollary}

%% ---------------------------------------------------------------------------
\section{Perception as Structural Matching}
\label{sec:perception}

The standard model of perception is:
\[
  \text{input} \;\xrightarrow{\;\text{processing}\;}\; \text{representation.}
\]
Under $H_0$ and Thesis~O, this model is replaced by a structural matching
problem.

\begin{definition}[Perception as SAT]
Perceiving for agent $A$ is defined as
\[
  \mathrm{Perceive}(A)
  = \operatorname*{arg\,max}_{S \subseteq \Univ} \CohTriple(S, \MA),
\]
where $\CohTriple(S, \MA)$ measures structural compatibility of $S$ with the
agent's model.
\end{definition}

This reformulation identifies perception with a constraint-satisfaction problem
(SAT) over the coherent sub-structures of the universe.  The agent does not
\emph{receive} a signal; it \emph{finds} the subset of $\Univ$ most compatible
with its current model.

\begin{remark}
This has a direct implication for semantic systems: document querying is not
retrieval but matching.  A query result is semantically valid if and only if
$\CohTriple(\text{result}) = (1,1,1)$.
\end{remark}

%% ---------------------------------------------------------------------------
\section{Lexical Coherence}
\label{sec:lexical}

\subsection{Definitions}

\begin{definition}[Structural Lexical Coherence, LC1]
A mapping $\phi : \mathcal{D}_1 \to \mathcal{D}_2$ between domains is
\emph{lexically coherent} if it preserves the coherence triple:
\[
  \forall S \subseteq \mathcal{D}_1,\quad
  \CohTriple(S) = (1,1,1)
  \;\Rightarrow\;
  \CohTriple(\phi(S)) = (1,1,1).
\]
A homomorphism satisfying this condition is called a \emph{coherence map}.
\end{definition}

\begin{definition}[Operational Lexical Coherence, LC2]
A coherence map $\phi$ is \emph{operationally valid} if it preserves the
retention score above a domain-specific threshold $\theta$:
\[
  \phi \text{ is valid} \;\iff\; \rho\bigl(\phi(S)\bigr) \geq \theta,
\]
where $\rho$ is the retention metric of \STF{}-\SIR{}.
\end{definition}

LC2 operationalizes LC1 for the \STF{}-\SIR{} pipeline: the existing retention
score $\rho$ is the computational proxy for coherence preservation under
domain transfer.

\subsection{Lexical Failure Taxonomy}

When $\phi$ fails to satisfy LC1 or LC2, the failure is classified into four
named modes.

\begin{center}
\renewcommand{\arraystretch}{1.25}
\begin{tabular}{llll}
\toprule
Tag & Name & Description & Error triple \\
\midrule
\texttt{LEX\_COLLAPSE} & Collapse &
  Source relations destroyed & $(0,{-},{-})$ \\
\texttt{LEX\_DRIFT} & Drift &
  Meaning deviates without detectable contradiction & $(1,0,0)$ \\
\texttt{LEX\_SPLIT} & Branching &
  One proposition maps to multiple interpretations & $(1,1,0)$ \\
\texttt{LEX\_MASK} & Mask &
  Mapping appears coherent but conceals incoherence & $(1,0,0)$ \\
\bottomrule
\end{tabular}
\end{center}

\texttt{LEX\_MASK} is the lexical analogue of hallucination: it is superficially
valid but structurally empty.

%% ---------------------------------------------------------------------------
\section{Coherence Error Taxonomy}
\label{sec:errors}

\subsection{Uniform Error Representation}

\begin{definition}[Error Triple]
The coherence error state of an item $x$ is the triple
\[
  \ErrTriple(x) = (\mathrm{Coh}_l,\; \mathrm{Coh}_g,\; \mathrm{Ground})
  \;\in\; \{0,1\}^3,
\]
where $\mathrm{Coh}_l$ is local coherence, $\mathrm{Coh}_g$ is global coherence,
and $\mathrm{Ground}$ is referential grounding.
\end{definition}

\subsection{Anomaly}

\begin{definition}[Anomaly, E1]
An item $x$ is an \emph{anomaly} with respect to domain $\mathcal{D}$ if
\[
  \mathrm{Anomaly}(x, \mathcal{D}) \;\iff\; P_{\mathcal{D}}(x) < \epsilon_{\mathcal{D}}.
\]
Error triple: $(1, 1, 1)$.
\end{definition}

An anomaly is statistically improbable but structurally valid.  It is recoverable
via Bayesian prior update.  Crucially, anomaly and a correct item share the same
error triple $(1,1,1)$; the distinction is probabilistic, not structural.

\subsection{Hallucination}

\begin{definition}[Hallucination, E2]
\label{def:hallucination}
An item $x$ is a \emph{hallucination} if
\[
  \mathrm{Hallucination}(x) \;\iff\;
  \mathrm{Coh}_l(x) = 1 \;\wedge\; \mathrm{Ground}(x, \mathcal{W}) = 0.
\]
Error triple: $(1, 0, 0)$.
\end{definition}

\begin{theorem}[Hallucination as Local Coherence Without Global Grounding]
\label{thm:hallucination}
A hallucination passes all internal coherence checks and is therefore
indistinguishable from a valid output by any purely logical verifier.  Detection
requires an external grounding mechanism.
\end{theorem}

\begin{proof}
By Definition~\ref{def:hallucination}, $\mathrm{Coh}_l(x) = 1$, so no internal
consistency check raises a flag.  The grounding failure $\mathrm{Ground}(x,
\mathcal{W}) = 0$ is only observable via comparison with an external referent.
No algorithm operating solely on $x$ can determine its grounding status.~$\square$
\end{proof}

Hallucination is constitutively distinct from lying: there is no intent; the
system operates outside its validity domain.

\subsection{Contradiction}

\begin{definition}[Contradiction, E3]
A set of propositions $S$ is \emph{contradictory} if
\[
  S \models \bot.
\]
Error triple: $(0, -, -)$.
\end{definition}

Contradiction fails at the local level, making global and grounding dimensions
irrelevant.  Unlike anomaly and hallucination, contradiction is formally
detectable without external data.

\subsection{Error Space Summary}

\begin{center}
\renewcommand{\arraystretch}{1.25}
\begin{tabular}{lcccc}
\toprule
Error type & $\mathrm{Coh}_l$ & $\mathrm{Coh}_g$ & Ground
           & Detection strategy \\
\midrule
Anomaly      & 1 & 1 & 1 & Corpus statistics over domain model \\
Hallucination & 1 & 0 & 0 & $\Delta$-tracking + external grounding \\
Contradiction & 0 & --- & --- & Formal verification / SAT \\
Valid        & 1 & 1 & 1 & --- \\
\bottomrule
\end{tabular}
\end{center}

%% ---------------------------------------------------------------------------
\section{Connection to STF-SIR}
\label{sec:stfsir}

\subsection{Evolved Token Model}

The \STF{}-\SIR{} v1 semantic token is
\[
  t = (L, S, \Sigma, \Gamma, C, P, \Delta, \Omega).
\]
This paper proposes two derived coordinates to be materialized during compilation:
\[
  t' = (L, S, \Sigma, \Gamma, C, P, \Delta, \Omega, \kappa, \iota),
\]
where
\[
  \kappa(t) = (C_l, C_c, C_o) \;\in\; \{0,1\}^3
\]
is the \emph{coherence coordinate}, and
\[
  \iota(t) = \bigl(H(t),\; \Inf(t),\; \ICE(t)\bigr)
  \;\in\; \mathbb{R}_{\geq 0}^2 \times \{0,1\}
\]
is the \emph{information coordinate}.

Both $\kappa$ and $\iota$ are \emph{derived} fields, never compiler inputs.
They are evaluated during the logical and semantic passes respectively, and
are available to the audit pipeline, diff engine, and retention benchmark.

A token satisfying $\kappa = (1,1,1)$ and $\ICE(t) = 1$ is a \emph{fully
coherent and productive semantic unit}.

\subsection{Hallucination in the Compilation Pipeline}

By Theorem~\ref{thm:hallucination}, a hallucinated token $t$ satisfies
$\omega_{\mathrm{inv}} = 1$ and $\omega_\rho \approx 1$ but
$\mathrm{Ground}(t, \mathcal{A}_{\mathrm{source}}) = 0$.  Detection therefore
requires provenance tracking in the $\Delta$ dimension.  Internal coherence
metrics are insufficient; a $\Delta$-tracking mechanism must be introduced in
the compilation pipeline to flag hallucinated tokens with the
\texttt{HALLUCINATION} tag.

\subsection{Query Engine Implication}

Under the perception-as-matching model of Section~\ref{sec:perception}, the
\STF{}-\SIR{} semantic query engine (EPIC-203) is a coherence-constrained SAT
solver over the SirGraph.  A query result is semantically valid if and only if
its coherence triple is $(1,1,1)$.

%% ---------------------------------------------------------------------------
\section{Implications}
\label{sec:implications}

\paragraph{For AI systems.}
Failures in large language models and generative systems can be reclassified as
coherence-triple mismatches.  Hallucinations (\S\ref{sec:errors}) are precisely
$(1,0,0)$ events: locally fluent outputs with zero grounding.  Grounding
mechanisms are, formally, devices for enforcing $\mathrm{Ground} = 1$.

\paragraph{For epistemology.}
Knowledge is bounded inference: an agent knows exactly what lies within its
inference horizon $\Horizon_A$.  Truth and coherence are orthogonal: a system
can be $(1,1,1)$ and false, or $(0,-,-)$ and accidentally true.

\paragraph{For systems design.}
A semantic compilation system is not a processor of tokens---it is a
\emph{coherence compiler}: a system whose primary invariant is the preservation
of the coherence triple across every transformation.  The retention score
$\rho \geq \theta$ is the operational expression of this invariant.

%% ---------------------------------------------------------------------------
\section{Conclusion}

We have shown that coherence, formalized as the triple $(C_l, C_c, C_o)$, is
logically prior to information and provides a uniform foundation for
understanding perception, inference, knowledge transfer, and error classification.
The key results are:

\begin{enumerate}
  \item Coherence is a primitive; information is derived (Thesis~O).
  \item Information is usable if and only if it is coherently executable
        (Theorem~\ref{thm:usability}).
  \item The inferential horizon of a finite agent coincides with $\mathbf{P}$
        (Theorem~\ref{thm:horizon}).
  \item Hallucination is local coherence without global grounding
        (Theorem~\ref{thm:hallucination}), and is undetectable by purely internal
        means.
  \item Lexical coherence generalizes the \STF{}-\SIR{} retention score $\rho$
        to cross-domain transfer.
\end{enumerate}

These results provide a formal foundation for the ongoing development of the
\STF{}-\SIR{} v2 platform, in which coherence is the central auditable invariant
of all semantic transformations.

%% ===========================================================================
\appendix

\section{Formal Foundations of Coherence and Information}
\label{app:foundations}

This appendix provides the logical and model-theoretic underpinnings for the
definitions and theorems in the main text.  The presentation follows the
tradition of Tarski~\cite{tarski1936} for semantic consequence and draws on
standard results in proof theory and computational complexity.

%% ---------------------------------------------------------------------------
\subsection{Preliminaries}
\label{app:prelim}

\paragraph{Language and structures.}
Let $\mathcal{L}$ be a first-order language, $\Univ$ the set of all
well-formed formulas (wffs) of $\mathcal{L}$, $\vdash$ the syntactic
derivability relation, and $\models$ the semantic consequence relation in the
sense of Tarski~\cite{tarski1936}.

\paragraph{Models.}
A model is a pair $\mathfrak{M} = (D, I)$ where $D$ is a non-empty domain and
$I$ an interpretation function.  We write $\mathfrak{M} \models \varphi$ when
$\mathfrak{M}$ satisfies $\varphi$.

\begin{definition}[Consistency, A1]
A set $S \subseteq \Univ$ is \emph{consistent} if and only if $S \nvdash \bot$.
\end{definition}

%% ---------------------------------------------------------------------------
\subsection{Coherence as a Generalization of Consistency}
\label{app:coherence-gen}

\paragraph{Logical coherence (A2).}
We identify logical coherence with syntactic consistency:
\[
  C_l(S) = 1 \;\iff\; S \nvdash \bot.
\]

\begin{theorem}[Compactness, A1]
\label{thm:compactness}
If $\mathcal{L}$ is a first-order language, then $S$ is consistent if and only
if every finite subset of $S$ is consistent:
\[
  S \nvdash \bot \;\iff\;
  \forall S_0 \subseteq S,\; |S_0| < \infty,\; S_0 \nvdash \bot.
\]
\end{theorem}

\begin{proof}
This is the classical Compactness Theorem for first-order logic, a consequence
of the completeness theorem proved by G\"{o}del~\cite{godel1930}.~$\square$
\end{proof}

\paragraph{Consequence for computation.}
Theorem~\ref{thm:compactness} is foundational for tractable coherence checking:
global coherence over a potentially infinite set can be reduced to checking
finite subsets, which makes local verification strategies sound.

%% ---------------------------------------------------------------------------
\subsection{Computational Coherence}
\label{app:comp-coh}

Define the coherence decision problem:
\[
  \mathrm{COH} = \bigl\{\, S \subseteq \Univ \mid S \nvdash \bot \,\bigr\}.
\]

\begin{theorem}[Complexity of Coherence, A2]
\label{thm:coh-conp}
For propositional $\mathcal{L}$, $\mathrm{COH} \in \mathrm{co\text{-}NP}$.
\end{theorem}

\begin{proof}
The complement $\overline{\mathrm{COH}}$ is the set of inconsistent
propositional sets.  Verifying inconsistency reduces to satisfiability of
$\neg\bigwedge S$: a nondeterministic algorithm guesses an assignment and
checks it in polynomial time.  Hence $\overline{\mathrm{COH}} \in \mathbf{NP}$
by the NP-completeness of SAT~\cite{cook1971}, and therefore
$\mathrm{COH} \in \mathrm{co\text{-}NP}$.~$\square$
\end{proof}

\begin{definition}[Computational Coherence — Theoretical Ideal, A3]
\label{def:compcoh}
\emph{(Theoretical motivation — non-normative for STF-SIR implementation.)}
$C_c(S) = 1$ if and only if there exists an algorithm $\alpha \in \mathbf{P}$
that verifies $S \nvdash \bot$.  In general propositional logic,
$C_c(S) = 1$ only for syntactically restricted fragments where coherence
checking is known to be in $\mathbf{P}$ (e.g., Horn clauses, 2-SAT).
\end{definition}

\begin{definition}[Computational Coherence — STF-SIR Operational Definition, A3-impl]
\label{def:compcoh-impl}
In STF-SIR, $C_c$ is approximated by a step-budget check:
\[
  C_c^{\mathit{impl}}(S,\, B) =
  \begin{cases}
    \mathit{Satisfied} & \text{if the coherence check uses } \leq B \text{ steps,} \\
    \mathit{Violated}  & \text{if it exhausts } B \text{ steps.}
  \end{cases}
\]
The default budget is $B = \mathtt{RECOMMENDED\_STEP\_BUDGET} = 1\,000\,000$.
This is a proxy for tractability; it makes no claim about $\mathbf{P}$ vs.\
$\mathbf{NP}$ membership of the problem instance.
\end{definition}

\begin{remark}
The distinction between $C_l$ and $C_c$ is not merely theoretical.  A set may
be logically coherent ($C_l = 1$) yet exhaust the step budget ($C_c = 0$) for
large or adversarially structured inputs.  The complexity-class analysis of
Theorem~\ref{thm:coh-conp} motivates \emph{why} a budget matters; it does not
define the operational behaviour of STF-SIR.
\end{remark}

%% ---------------------------------------------------------------------------
\subsection{Operational Coherence}
\label{app:op-coh}

Let $\mathrm{Cn}(S) = \{p \in \Univ \mid S \vdash p\}$ denote the deductive
closure of $S$.

\begin{definition}[Operational Coherence, A4]
$C_o(S) = 1$ if and only if there exists $p \in \Univ$ such that
$S \vdash p$ and $p \notin \mathrm{Cn}(\emptyset)$, i.e., $S$ derives at
least one non-tautological consequence.
\end{definition}

\begin{theorem}[Operational Coherence via Models, A3]
\label{thm:op-sat}
If $S$ is finite and consistent, then $C_o(S) = 1$ if and only if there exists
a model $\mathfrak{M}$ such that $\mathfrak{M} \models S$.
\end{theorem}

\begin{proof}
($\Rightarrow$) Any non-tautological consequence $p$ derived from $S$ is true
in every model of $S$; hence $S$ has a model.

($\Leftarrow$) By G\"{o}del's Completeness Theorem~\cite{godel1930}: if $S$
has a model, then $S$ is consistent, and consistency of a non-trivial set
implies the existence of non-tautological derivable consequences.~$\square$
\end{proof}

%% ---------------------------------------------------------------------------
\subsection{Coherently Executable Information: Formal Proofs}
\label{app:cei-proofs}

Let an agent be a pair $A = (\MA, F_A)$ where $\MA \subseteq \Univ$ is its
theory and $F_A$ is its inference system.  The CEI definition from the main
text (Definition~\ref{def:cei}) is:
\[
  \ICE(m, A) = 1 \;\iff\;
  \begin{cases}
    \MA \cup \{m\} \nvdash \bot & \text{(C1)} \\
    \exists\, p \notin \mathrm{Cn}(\MA) : \MA \cup \{m\} \vdash p & \text{(C2)}
  \end{cases}
\]

\begin{theorem}[Necessity of Coherence, A4]
\label{thm:explosion}
If $\MA \cup \{m\} \vdash \bot$, then $\MA \cup \{m\} \vdash p$ for every
$p \in \Univ$.
\end{theorem}

\begin{proof}
This is the principle of explosion (\emph{ex contradictione quodlibet}):
$\bot \vdash p$ is a theorem of any standard logical calculus.~$\square$
\end{proof}

\paragraph{Consequence.}
By Theorem~\ref{thm:explosion}, an incoherent message does not merely fail to
add information---it makes every proposition derivable, destroying the
discriminative power of the agent's model.  This is why C1 is not merely
desirable but necessary.

\begin{theorem}[Inutility of Non-Execution, A5]
\label{thm:inutility}
If $\MA \cup \{m\} \nvdash p$ for every $p \notin \mathrm{Cn}(\MA)$, then $m$
adds no useful information to $A$.
\end{theorem}

\begin{proof}
If no new non-tautological consequence follows from $\MA \cup \{m\}$, then
$\mathrm{Cn}(\MA \cup \{m\}) \subseteq \mathrm{Cn}(\MA)$.  The message does
not expand the agent's knowledge; it is informationally inert.~$\square$
\end{proof}

%% ---------------------------------------------------------------------------
\subsection{Computational Limits of Cognition: Proofs}
\label{app:limits-proofs}

\begin{theorem}[Inference Horizon Bound, A6]
\label{thm:ha-p}
$\Horizon_A \subseteq \mathbf{P}$.
\end{theorem}

\begin{proof}
By definition, $p \in \Horizon_A$ requires the existence of $\alpha \in
\mathbf{P}$ such that $\alpha(\MA, O_A) \vdash p$.  Therefore every element
of $\Horizon_A$ is witnessed by a polynomial-time derivation.~$\square$
\end{proof}

\begin{theorem}[NP Barrier, A7]
\label{thm:np-barrier}
If $\mathbf{P} \neq \mathbf{NP}$, then there exist propositions $p$ such that
$p \notin \Horizon_A$ for every polynomial-time agent $A$.
\end{theorem}

\begin{proof}
Assuming $\mathbf{P} \neq \mathbf{NP}$, there exist decision problems in
$\mathbf{NP} \setminus \mathbf{P}$.  Any proposition whose derivation requires
solving such a problem cannot be produced by any $\alpha \in \mathbf{P}$.  By
definition of $\Horizon_A$, such propositions lie outside every agent's
inference horizon.~$\square$
\end{proof}

%% ---------------------------------------------------------------------------
\subsection{Perception as Satisfiability: Complexity}
\label{app:sat-complexity}

\begin{definition}[Perception as SAT, A8]
For agent $A$ with model $\MA$ and observation set $O$:
\[
  \mathrm{Perceive}(A, O) = \mathrm{SAT}(\MA \cup O).
\]
\end{definition}

\begin{theorem}[NP-Completeness of Perception, A8]
\label{thm:perception-np}
In the propositional case, $\mathrm{Perceive}(A, O)$ is NP-complete.
\end{theorem}

\begin{proof}
Membership in NP: a nondeterministic algorithm guesses a satisfying assignment
and verifies it in polynomial time.  NP-hardness: propositional SAT reduces to
$\mathrm{Perceive}$ by setting $\MA = \emptyset$ and $O$ to the formula in
question~\cite{cook1971}.~$\square$
\end{proof}

\begin{remark}
Theorem~\ref{thm:perception-np} formalizes the cognitive cost of alignment:
perception, as structural matching, is computationally harder than recognition
(polynomial) but tractable in restricted fragments.  Natural cognitive systems
likely exploit domain-specific structure to reduce perception to tractable
subcases.
\end{remark}

%% ---------------------------------------------------------------------------
\subsection{Lexical Coherence: Homomorphism Property}
\label{app:lexical-proofs}

\begin{theorem}[Homomorphism Preserves Coherence, A9]
\label{thm:hom-coherence}
If $\phi : \mathcal{L}_1 \to \mathcal{L}_2$ is a logical homomorphism
(preserving the inference rules of $\mathcal{L}_1$), then:
\[
  S \vdash p \;\Rightarrow\; \phi(S) \vdash \phi(p).
\]
Consequently, if $S$ is consistent in $\mathcal{L}_1$, then $\phi(S)$ is
consistent in $\mathcal{L}_2$.
\end{theorem}

\begin{proof}
Let $S \vdash p$ via a derivation $d = (f_1, \ldots, f_n)$ where $f_n = p$ and
each $f_i$ is either in $S$ or follows from earlier steps by a rule of
$\mathcal{L}_1$.  Since $\phi$ preserves the inference rules, applying $\phi$
pointwise to $d$ yields a valid derivation of $\phi(p)$ from $\phi(S)$ in
$\mathcal{L}_2$.  Taking $p = \bot$ contrapositively: if $\phi(S) \vdash \bot$,
then $S \vdash \bot$, so $C_l(S) = 0$.  The contrapositive of this is LC1.~$\square$
\end{proof}

%% ---------------------------------------------------------------------------
\subsection{Error Taxonomy: Formal Characterizations}
\label{app:error-proofs}

\begin{definition}[Hallucination via Satisfiability, A9]
An item $x$ over theory $S$ is a hallucination if $C_l(S) = 1$ (locally
consistent) and there is no model $\mathfrak{M}$ such that $\mathfrak{M}
\models S \cup \{\mathrm{context}(x)\}$ where $\mathrm{context}(x)$ includes
the grounding conditions for $x$.
\end{definition}

\begin{theorem}[Hallucination Implies Global Non-Satisfiability, A10]
\label{thm:hall-unsat}
A hallucination is locally coherent but globally unsatisfiable: the source
context $\mathcal{A}$ provides no model for the grounding conditions of $x$.
\end{theorem}

\begin{proof}
By definition, $C_l(x) = 1$ means $x$ does not produce a contradiction in
isolation.  The grounding failure $\mathrm{Ground}(x, \mathcal{W}) = 0$
means $\neg\exists\, \mathfrak{M} \in \mathcal{W} : \mathfrak{M} \models x$.
Therefore $x$ is satisfiable as a formula but has no model in the intended
world $\mathcal{W}$: local coherence and global non-satisfiability coexist.~$\square$
\end{proof}

%% ---------------------------------------------------------------------------
\subsection{Central Theorem}
\label{app:central}

We now state and prove the main result of the framework, consolidating
Theorems~\ref{thm:explosion}, \ref{thm:inutility}, and the definition of CEI.

\begin{theorem}[Primacy of Coherence: Central Result]
\label{thm:central}
Information is useful if and only if it is coherent and executable:
\[
  m \text{ is useful for } A
  \;\iff\;
  \ICE(m, A) = 1
  \;\iff\;
  C_l(\MA \cup \{m\}) = 1 \;\wedge\; C_o(\MA \cup \{m\}) = 1.
\]
\end{theorem}

\begin{proof}
($\Rightarrow$) Suppose $m$ is useful, i.e., it expands $A$'s knowledge
without destroying it.  If $\MA \cup \{m\} \vdash \bot$, then by
Theorem~\ref{thm:explosion} every proposition is derivable, and no information
is gained selectively---a contradiction with usefulness.  If $C_o = 0$, then
no new non-tautological consequence follows, again contradicting usefulness.
Therefore $C_l = 1$ and $C_o = 1$.

($\Leftarrow$) If $C_l(\MA \cup \{m\}) = 1$, the model is not destroyed.  If
$C_o(\MA \cup \{m\}) = 1$, at least one new non-tautological consequence
follows.  Both conditions jointly constitute the definition of usefulness.~$\square$
\end{proof}

\paragraph{Summary.}
The central theorem establishes the full hierarchy:
\begin{itemize}
  \item Coherence ($C_l$) is the necessary condition: without it, information
        is destructive.
  \item Executability ($C_o$) is the sufficient productive condition: without
        it, information is inert.
  \item Truth is orthogonal: a useful message can be false (it expands the
        model coherently but incorrectly), and a true message can be useless
        (tautological).
\end{itemize}
This completes the formal grounding of Thesis~O: coherence is logically and
computationally prior to information.

%% ---------------------------------------------------------------------------
\begin{thebibliography}{9}

\bibitem{shannon1948}
  C.~E. Shannon,
  ``A mathematical theory of communication,''
  \textit{Bell System Technical Journal}, vol.~27, pp.~379--423 and 623--656,
  1948.

\bibitem{carnap1953}
  R.~Carnap and Y.~Bar-Hillel,
  ``An outline of a theory of semantic information,''
  Technical Report No.~247, MIT Research Laboratory of Electronics, 1952.
  Published as: \textit{Language and Information}, Addison-Wesley, 1964.

\bibitem{minkowski1908}
  H.~Minkowski,
  ``Raum und Zeit,''
  Address delivered at the 80th Assembly of German Natural Scientists and
  Physicians, Cologne, 1908.

\bibitem{cherniak1986}
  C.~Cherniak,
  \textit{Minimal Rationality}.
  MIT Press, Cambridge, MA, 1986.

\bibitem{ji2023}
  Z.~Ji et al.,
  ``Survey of hallucination in natural language generation,''
  \textit{ACM Computing Surveys}, vol.~55, no.~12, pp.~1--38, 2023.

\bibitem{figueiredo2026sts}
  R.~Figueiredo,
  ``Semantic Token Space (\STS{}): A Multidimensional Formalization for
  \STF{}-\SIR{},''
  Technical report, 2026.
  Available: \texttt{docs/sts-formalization.md}.

\bibitem{figueiredo2026spec}
  R.~Figueiredo,
  ``\STF{} Core Semantics Specification,''
  Technical report, 2026.
  Available: \texttt{docs/spec/SPEC-STF-CORE-SEMANTICS.md}.

\bibitem{tarski1936}
  A.~Tarski,
  ``\"{U}ber den Begriff der logischen Folgerung,''
  \textit{Actes du Congr\`{e}s International de Philosophie Scientifique},
  vol.~7, pp.~1--11, 1936.
  English translation: ``On the concept of logical consequence,''
  in \textit{Logic, Semantics, Metamathematics}, Oxford University Press, 1956.

\bibitem{godel1930}
  K.~G\"{o}del,
  ``Die Vollst\"{a}ndigkeit der Axiome des logischen Funktionenkalk\"{u}ls,''
  \textit{Monatshefte f\"{u}r Mathematik und Physik}, vol.~37, pp.~349--360,
  1930.
  English translation: ``The completeness of the axioms of the functional
  calculus of logic,'' in \textit{From Frege to G\"{o}del}, Harvard University
  Press, 1967.

\bibitem{cook1971}
  S.~A. Cook,
  ``The complexity of theorem-proving procedures,''
  in \textit{Proceedings of the 3rd Annual ACM Symposium on Theory of
  Computing (STOC)}, pp.~151--158, 1971.

\end{thebibliography}

\end{document}
